% SPDX-License-Identifier: CC-BY-4.0
\documentclass[11pt,letterpaper]{article}

\usepackage[T1]{fontenc}
\usepackage{amsmath,amssymb}
\usepackage{tikz}
\usepackage{pgfplots}
\pgfplotsset{compat=1.18}
\usepackage{algorithm}
\usepackage{algpseudocode}
\usepackage{booktabs}
\usepackage{enumitem}
\usepackage[hidelinks]{hyperref}

\newcommand{\N}{\mathbb{N}}
\newcommand{\Z}{\mathbb{Z}}
\newcommand{\Q}{\mathbb{Q}}
\newcommand{\R}{\mathbb{R}}
\newcommand{\C}{\mathbb{C}}
\newcommand{\Map}[3]{#1 \colon #2 \to #3}
\newcommand{\Mod}[1]{#1\text{-}\mathrm{Mod}}
\newcommand{\GroupHom}{\mathrm{Hom}_{\mathbf{Grp}}}
\newcommand{\RingHom}{\mathrm{Hom}_{\mathbf{Ring}}}

\newtheorem{theorem}{Theorem}
\newtheorem{lemma}[theorem]{Lemma}
\newtheorem{corollary}[theorem]{Corollary}
\newtheorem{definition}[theorem]{Definition}
\newtheorem{remark}[theorem]{Remark}
\newenvironment{proof}{\noindent\textit{Proof.}}{\hfill\(\square\)}

\title{Generalized Collatz Library}
\author{}
\date{\today}

\begin{document}
\maketitle

\begin{abstract}
We classify the positive solutions of \(x^2+D=2^n\) for \(D=7\) and \(D=23\) at every even exponent and at all exponents \(n\le 32\). A difference-of-squares identity also gives two distinct positive representations for every Mersenne number \(D=2^k-1\) with \(k\ge 4\). The bounded classifications are supported by explicit integer-square brackets and matching Lean proofs.
\end{abstract}

\section{The displayed pairs}

Throughout, \(D\), \(n\), and \(x\) are natural numbers, and a positive solution is a pair with \(x>0\).

\begin{theorem}
\label{thm:displayed}
The following seven identities hold:
\begin{align*}
1^2+7&=2^3,\\
3^2+7&=2^4,\\
5^2+7&=2^5,\\
11^2+7&=2^7,\\
181^2+7&=2^{15},\\
3^2+23&=2^5,\\
45^2+23&=2^{11}.
\end{align*}
\end{theorem}

\begin{proof}
The first four evaluations are \(1+7=8\), \(9+7=16\), \(25+7=32\), and \(121+7=128\). For the fifth, \(180^2=32400\) and
\[
181^2=180^2+2\cdot 180+1=32761,
\]
so \(181^2+7=32768=2^{15}\). The last two evaluations are \(9+23=32\) and \(45^2=2025\), hence \(2025+23=2048=2^{11}\).
\end{proof}

\section{A square gap}

\begin{lemma}
\label{lem:gap}
For every natural number \(m\),
\[
(2^m)^2-(2^m-1)^2=2^{m+1}-1.
\]
In particular \((2^m)^2=2^{2m}\).
\end{lemma}

\begin{proof}
The difference of squares is
\[
(2^m)^2-(2^m-1)^2=\bigl(2^m+(2^m-1)\bigr)\bigl(2^m-(2^m-1)\bigr).
\]
The second factor equals \(1\). Since \(2^{m+1}=2^m+2^m\),
\[
2^m+(2^m-1)+1=2^{m+1},
\]
and therefore \(2^m+(2^m-1)=2^{m+1}-1\). The product is \(2^{m+1}-1\). The power identity \((2^m)^2=2^{m\cdot 2}=2^{2m}\) is the rule \(a^{uv}=(a^u)^v\).
\end{proof}

\begin{lemma}
\label{lem:between}
Let \(a\) and \(t\) be natural numbers with \(a^2<t<(a+1)^2\). No natural number \(x\) satisfies \(x^2=t\).
\end{lemma}

\begin{proof}
If \(x\le a\), then \(x^2\le a^2<t\). If \(x\ge a+1\), then \(x^2\ge(a+1)^2>t\).
\end{proof}

\begin{lemma}
\label{lem:strict-seven}
If \(m\ge 3\), then
\[
(2^m-1)^2<2^{2m}-7<(2^m)^2.
\]
\end{lemma}

\begin{proof}
Lemma \ref{lem:gap} gives
\[
(2^m-1)^2+(2^{m+1}-1)=(2^m)^2.
\]
Since \(m\ge 3\), one has \(m+1\ge 4\) and \(2^{m+1}\ge 2^4=16\). Thus \(8<2^{m+1}\), so \(7<2^{m+1}-1\). Adding \((2^m-1)^2\) yields
\[
(2^m-1)^2+7<(2^m)^2.
\]
Subtracting \(7\) gives the lower bound \((2^m-1)^2<2^{2m}-7\). For the upper bound, \(2m\ge 6>4\), so \(2^{2m}\ge 16>7\) and \(2^{2m}-7<2^{2m}=(2^m)^2\).
\end{proof}

\begin{lemma}
\label{lem:strict-twenty}
If \(m\ge 4\), then
\[
(2^m-1)^2<2^{2m}-23<(2^m)^2.
\]
\end{lemma}

\begin{proof}
The same expansion as in Lemma \ref{lem:strict-seven} applies. Here \(m+1\ge 5\), so \(2^{m+1}\ge 32\). Hence \(24<2^{m+1}\) and \(23<2^{m+1}-1\). Therefore
\[
(2^m-1)^2+23<(2^m)^2,
\]
which is the lower bound. Also \(2m\ge 8\), so \(2^{2m}\ge 256>23\) and the upper bound follows.
\end{proof}

\section{Even exponents}

\begin{theorem}
\label{thm:even-seven}
For a positive integer \(x\) and a natural number \(m\),
\[
x^2+7=2^{2m}
\]
if and only if \(x=3\) and \(m=2\).
\end{theorem}

\begin{proof}
The identity \(3^2+7=16=2^4\) is Theorem \ref{thm:displayed}, and \(4=2\cdot 2\).

Conversely, suppose \(x>0\) and \(x^2+7=2^{2m}\). If \(m=0\), then \(x^2+7=1\), which is impossible because the left side is at least \(7\). If \(m=1\), then \(x^2+7=4\), which is likewise impossible. If \(m=2\), then \(x^2=9\). The bound \(x\ge 4\) would give \(x^2\ge 16>9\), and the bound \(x\le 2\) would give \(x^2\le 4<9\), so \(x=3\).

If \(m\ge 3\), Lemma \ref{lem:strict-seven} places \(2^{2m}-7\) strictly between \((2^m-1)^2\) and \((2^m)^2\). Lemma \ref{lem:between} says that this remainder is not a square, contradicting \(x^2=2^{2m}-7\).
\end{proof}

\begin{theorem}
\label{thm:even-twenty}
For every natural number \(m\) and every positive integer \(x\), \(x^2+23\neq 2^{2m}\).
\end{theorem}

\begin{proof}
If \(m\le 2\), then \(2^{2m}\le 16<23\le x^2+23\). If \(m=3\), then \(x^2=2^6-23=41\). Since \(6^2=36<41<49=7^2\), Lemma \ref{lem:between} applies. If \(m\ge 4\), Lemma \ref{lem:strict-twenty} places \(2^{2m}-23\) strictly between \((2^m-1)^2\) and \((2^m)^2\), so again the remainder is not a square.
\end{proof}

\section{The odd box of height \(32\)}

\begin{lemma}
\label{lem:root}
Let \(D\) and \(n\) be natural numbers and let \(s\) be a positive integer. Then \(s^2+D=2^n\) if and only if \(2^n\ge D\) and \(s^2=2^n-D\). If \(a^2\le 2^n-D<(a+1)^2\), the remainder is a square if and only if it equals \(a^2\).
\end{lemma}

\begin{proof}
If \(s^2+D=2^n\) and \(s>0\), then \(D\le 2^n\) and \(s^2=2^n-D\). The converse adds \(D\). For the bracket assertion, suppose \(r=2^n-D=x^2\). If \(x<a\), then \(x^2<a^2\le r\), a contradiction. If \(x>a\), then \(x\ge a+1\) and \(x^2\ge(a+1)^2>r\), again a contradiction. Thus \(x=a\). Conversely, \(r=a^2\) is a square.
\end{proof}

\begin{theorem}
\label{thm:box-seven}
Let \(x\) be a positive integer and let \(n\le 32\). Then \(x^2+7=2^n\) if and only if \((x,n)\) is one of
\[
(1,3),\ (3,4),\ (5,5),\ (11,7),\ (181,15).
\]
\end{theorem}

\begin{proof}
Each listed pair satisfies the equation by Theorem \ref{thm:displayed}. For the converse, write \(n=2m\) or \(n=2m+1\).

If \(n\) is even, Theorem \ref{thm:even-seven} forces \((x,n)=(3,4)\).

If \(n\) is odd and \(n\le 32\), Lemma \ref{lem:root} asks whether \(2^n-7\) is a positive square, except for \(n=1\), where \(2<7\) and no solution exists. The odd values from \(3\) through \(31\) have the following brackets. Equality is a solution; a strict bracket is not.

\begin{center}
\begin{tabular}{rrrrr}
\toprule
\(n\) & \(2^n-7\) & \(s\) & \(s^2\) & \((s+1)^2\) \\
\midrule
\(3\) & \(1\) & \(1\) & \(1\) & \(4\) \\
\(5\) & \(25\) & \(5\) & \(25\) & \(36\) \\
\(7\) & \(121\) & \(11\) & \(121\) & \(144\) \\
\(9\) & \(505\) & \(22\) & \(484\) & \(529\) \\
\(11\) & \(2041\) & \(45\) & \(2025\) & \(2116\) \\
\(13\) & \(8185\) & \(90\) & \(8100\) & \(8281\) \\
\(15\) & \(32761\) & \(181\) & \(32761\) & \(33124\) \\
\(17\) & \(131065\) & \(362\) & \(131044\) & \(131769\) \\
\(19\) & \(524281\) & \(724\) & \(524176\) & \(525625\) \\
\(21\) & \(2097145\) & \(1448\) & \(2096704\) & \(2099601\) \\
\(23\) & \(8388601\) & \(2896\) & \(8386816\) & \(8392609\) \\
\(25\) & \(33554425\) & \(5792\) & \(33547264\) & \(33558849\) \\
\(27\) & \(134217721\) & \(11585\) & \(134212225\) & \(134235396\) \\
\(29\) & \(536870905\) & \(23170\) & \(536848900\) & \(536895241\) \\
\(31\) & \(2147483641\) & \(46340\) & \(2147395600\) & \(2147488281\) \\
\bottomrule
\end{tabular}
\end{center}

The rows \(n=3,5,7,15\) are equalities and recover \((1,3)\), \((5,5)\), \((11,7)\), and \((181,15)\). In every other row, \(s^2<2^n-7<(s+1)^2\), so Lemma \ref{lem:between} excludes a solution.
\end{proof}

\begin{theorem}
\label{thm:box-twenty}
Let \(x\) be a positive integer and let \(n\le 32\). Then \(x^2+23=2^n\) if and only if \((x,n)\) is one of \((3,5)\) and \((45,11)\).
\end{theorem}

\begin{proof}
Both pairs occur in Theorem \ref{thm:displayed}. If \(n\) is even, Theorem \ref{thm:even-twenty} forbids a positive solution. If \(n\) is odd and \(n\le 3\), then \(2^n<23\). The remaining odd exponents have the brackets

\begin{center}
\begin{tabular}{rrrrr}
\toprule
\(n\) & \(2^n-23\) & \(s\) & \(s^2\) & \((s+1)^2\) \\
\midrule
\(5\) & \(9\) & \(3\) & \(9\) & \(16\) \\
\(7\) & \(105\) & \(10\) & \(100\) & \(121\) \\
\(9\) & \(489\) & \(22\) & \(484\) & \(529\) \\
\(11\) & \(2025\) & \(45\) & \(2025\) & \(2116\) \\
\(13\) & \(8169\) & \(90\) & \(8100\) & \(8281\) \\
\(15\) & \(32745\) & \(180\) & \(32400\) & \(32761\) \\
\(17\) & \(131049\) & \(362\) & \(131044\) & \(131769\) \\
\(19\) & \(524265\) & \(724\) & \(524176\) & \(525625\) \\
\(21\) & \(2097129\) & \(1448\) & \(2096704\) & \(2099601\) \\
\(23\) & \(8388585\) & \(2896\) & \(8386816\) & \(8392609\) \\
\(25\) & \(33554409\) & \(5792\) & \(33547264\) & \(33558849\) \\
\(27\) & \(134217705\) & \(11585\) & \(134212225\) & \(134235396\) \\
\(29\) & \(536870889\) & \(23170\) & \(536848900\) & \(536895241\) \\
\(31\) & \(2147483625\) & \(46340\) & \(2147395600\) & \(2147488281\) \\
\bottomrule
\end{tabular}
\end{center}

Equality holds only for \(n=5\) and \(n=11\). Every other row is a strict bracket.
\end{proof}

\section{Mersenne representations}

\begin{theorem}
\label{thm:mersenne}
For every integer \(k\ge 4\),
\[
1^2+(2^k-1)=2^k,
\]
\[
\bigl(2^{k-1}-1\bigr)^2+(2^k-1)=2^{2k-2},
\]
the two positive roots are distinct, and the two exponents are distinct.
\end{theorem}

\begin{proof}
The first identity is \(1+(2^k-1)=2^k\). For the second, set \(m=k-1\). Then \(m\ge 3\), \(k=m+1\), and \(2k-2=2m\). Lemma \ref{lem:gap} rearranges to
\[
(2^m-1)^2+(2^{m+1}-1)=(2^m)^2=2^{2m}.
\]
Substituting \(m=k-1\) is the second identity. The roots differ because \(k-1\ge 3\), so \(2^{k-1}-1\ge 2^3-1=7>1\). The exponents differ because \(2k-2-k=k-2\ge 2\).
\end{proof}

\begin{remark}
\label{rem:not}
The even-exponent conclusions and the Mersenne identities are uniform. The odd-exponent classifications apply only through \(n=32\). The Mersenne identities do not classify all parameters \(D\) admitting two representations.
\end{remark}

\end{document}
